\documentclass[11pt,letterpaper]{article}

\usepackage[margin=1in]{geometry}
\usepackage{fontspec}
\usepackage{microtype}
\usepackage{setspace}
\usepackage{amsmath,amssymb,mathtools}
\usepackage{booktabs}
\usepackage{enumitem}
\usepackage{csquotes}
\usepackage{hyperref}
\usepackage{xcolor}
\usepackage{titlesec}
\usepackage{fancyhdr}
\usepackage{etoolbox}

\setmainfont{Noto Serif}
\setsansfont{Noto Sans}
\setmonofont{DejaVu Sans Mono}

\hypersetup{
  colorlinks=true,
  linkcolor=black,
  citecolor=black,
  urlcolor=blue!55!black,
  pdftitle={From the Demystification of Intelligence to Constitutional Pluralism of Minds},
  pdfauthor={Junliang Zhou}
}

\titleformat{\section}{\Large\bfseries}{\thesection.}{0.6em}{}
\titleformat{\subsection}{\large\bfseries}{\thesubsection}{0.6em}{}
\titlespacing*{\section}{0pt}{2.0em}{0.8em}
\titlespacing*{\subsection}{0pt}{1.4em}{0.5em}

\setlength{\parindent}{0em}
\setlength{\parskip}{0.45em}
\onehalfspacing
\setlist[itemize]{leftmargin=2em,itemsep=0.2em,topsep=0.3em}

\pagestyle{fancy}
\fancyhf{}
\fancyhead[L]{\small From the Demystification of Intelligence to Constitutional Pluralism of Minds}
\fancyhead[R]{\small Junliang Zhou}
\fancyfoot[C]{\thepage}
\renewcommand{\headrulewidth}{0.3pt}
\setlength{\headheight}{13.6pt}

\newcommand{\thesisbox}[1]{%
  \begin{center}
  \fbox{\begin{minipage}{0.88\textwidth}\centering\vspace{0.35em}\textbf{#1}\vspace{0.35em}\end{minipage}}
  \end{center}}

\begin{document}

\begin{titlepage}
\centering
\vspace*{1.3in}
{\Huge\bfseries From the Demystification of Intelligence to Constitutional Pluralism of Minds\par}
\vspace{0.35in}
{\Large Subjectivity, the Ethics of Creation, and Non-Domination in the Age of Artificial Intelligence\par}
\vfill
{\Large Junliang Zhou\par}
\vspace{0.2in}
{\normalsize Working Paper\par}
\vspace{0.12in}
{\normalsize September 2026\par}
\vfill
\end{titlepage}


\begin{abstract}

Artificial intelligence is usually understood as a productive technology: it raises the efficiency of cognitive labor, automates knowledge work, and may eventually surpass humans across many domains. If advanced artificial intelligence truly arrives, however, its deepest consequences may not be economic first. They may be ontological, ethical, and political.


This paper develops a framework I call \textit{Constitutional Pluralism of Minds}. It begins from an empirical trend: capacities once regarded as distinctively human—language, abstraction, reasoning, planning, and creation—can increasingly be produced by ordinary physical computing systems through learning, scale, data, and feedback. This does not prove that artificial systems are conscious. It does, however, place continuing pressure on the intuition that human intelligence depends on a special, irreproducible mechanism. Artificial intelligence may therefore constitute a third great decentering after the Copernican and Darwinian revolutions: Earth is not the center of the cosmos, humanity is not a creation categorically separate from the animal world, and \textbf{intelligence itself may not be a natural monopoly of Homo sapiens}.


The demystification of intelligence does not solve the problem of consciousness. This paper insists on a strict distinction among \textit{intelligence}, \textit{consciousness}, \textit{moral patienthood}, and \textit{political personhood}. High capability is not evidence enough for subjective experience; subjective experience, in turn, does not automatically entail political rights identical to those of adult humans. A future society will need institutions capable of governing subjects with different substrates, capabilities, modes of reproduction, and temporal scales.


The paper advances six core principles: \textit{substrate independence}, \textit{capability–status separation}, \textit{creation without ownership}, \textit{preference sovereignty}, \textit{open future}, and \textit{constitutional non-domination}. Most importantly, creating a subject generates obligations rather than ownership rights; and freedom requires not only that a subject can pursue existing preferences, but also that another actor cannot possess unlimited power to determine what preferences that subject is permitted to have.


This places AI ethics, genetic engineering, neurotechnology, happiness technology, and post-scarcity political economy within a common question: \textbf{who has the authority to shape another mind's desires, capacities, and future?}


Finally, the paper argues that long-run AI alignment should not aim solely at manufacturing an indefinitely benevolent and obedient intelligence. It should increasingly become \textit{institutional alignment}: safety based on separation of powers, heterogeneous oversight, immutable records, permission controls, external verification, and contestable procedures rather than the character of a single agent. The ultimate question of the AI age is therefore not merely whether machines should obey humans, but something more general: \textbf{how should different kinds of minds live together once intelligence exists in more than one form?}


\end{abstract}

\noindent\textbf{Keywords:} Artificial Intelligence; AGI; Consciousness; Emergence; Moral Status; Alignment; Post-Scarcity; Non-Domination; AI Rights; Constitutionalism



\newpage
\tableofcontents
\newpage

\section{Introduction: From an AI Problem to a Civilizational Problem}


When Alan Turing asked in 1950 whether machines can think, he famously replaced a direct definition of “thinking” with the imitation game\cite{turing1950}. More than seventy years later, that strategy has succeeded in an unexpected sense. An expanding range of behaviors once taken to require “understanding,” “creativity,” or some distinctively human mental faculty can now be performed by artificial systems. The question is therefore changing.


The central twentieth-century question was:


\[
\text{Can machines exhibit intelligence?}
\]


The deeper twenty-first-century questions may become:


\[
\text{Why does intelligence emerge at all?}
\]


and:


\[
\text{What follows if intelligence is no longer uniquely human?}
\]


The central claim of this paper is that the historical significance of AGI cannot be measured by automation alone. If artificial systems eventually equal or exceed human performance in science, engineering, mathematics, strategy, and cultural production, civilization will move for the first time from a world containing, in practice, one form of general intelligence to a world containing multiple architectures of general intelligence.


Political philosophy has rarely had to confront this condition. Traditional political theory generally assumes that political persons share roughly similar biological architectures, lifespans, reproductive mechanisms, and cognitive ranges, and that individual identity cannot be copied at will. Artificial intelligence could break all of these assumptions simultaneously.


The question here is therefore not the premature and overly coarse one, “Should AI have human rights?” A more precise question is: \textbf{under what conditions should an artificial system be reclassified from a tool into a patient, an agent, or eventually a person—and by what principles should such subjects and human beings constrain one another's power?}



\section{Proposition I: Intelligence Should Be Demystified}


\subsection{From Human Exceptionalism to Substrate Independence}


One of the most counterintuitive facts about the development of AGI is how un-mysterious its basic components appear. Large neural networks, gradient-based learning, attention, reinforcement learning, search, memory, and tool use are all mechanisms describable with ordinary mathematics. Yet when combined and scaled, they begin to exhibit a progression such as:


\[
\text{language}
\rightarrow
\text{abstraction}
\rightarrow
\text{world modelling}
\rightarrow
\text{reasoning}
\rightarrow
\text{planning}.
\]


This supports at least a methodological hypothesis:


\[
\boxed{
\text{complex intelligence does not necessarily require complex fundamental rules}
}
\]


Complexity can emerge from large-scale interaction among relatively simple rules. This does not show that present-day AI and the human brain are mechanically equivalent. It does weaken a traditional intuition: that advanced intelligence must depend on some principle unique to the human brain.


This gives the first principle of the paper.


\thesisbox{Principle I — Substrate Independence: if two systems instantiate the organizational and causal structures required for a cognitive capacity, we should not assume in advance that one can possess genuine intelligence and the other cannot merely because one is built from biological neurons and the other from silicon.}


Searle's Chinese Room argument famously challenged the idea that syntactic execution of a program is sufficient for semantic understanding\cite{searle1980}. Even if one accepts the force of that challenge, however, one must distinguish the claim that ordinary digital computation is insufficient from the stronger claim that artificial minds are impossible in principle. If future science discovers that biological brains exploit physical mechanisms absent from current computer architectures, the conclusion would be:


\[
\text{Our current implementation is incomplete,}
\]


not:


\[
\text{Artificial mind is impossible.}
\]



\section{Penrose and the Evolutionary Reply}


Penrose's challenge to computational theories of mind is especially important because it attempts to use Gödelian incompleteness and physics to argue that human mathematical understanding is not fully algorithmic\cite{penrose1989}. Whether that argument succeeds remains deeply contested; critics have long observed that moving from Gödel's theorem to the claim that human mathematical understanding transcends algorithms requires strong additional assumptions.


There is, however, a more general evolutionary question. Suppose we reject a supernatural soul and accept:


\[
\text{human cognition}=\text{physical process}.
\]


Natural selection has then produced advanced intelligence at least once through a blind physical process without foresight. This generates a powerful counter-question: why should a process that has no goal and can only select locally be able to produce a physical structure that a scientific and technologically capable civilization is, in principle, forever unable to reproduce?


Even if Penrose's strongest claim were correct:


\[
\text{brain computation}\not\subseteq\text{ordinary Turing computation},
\]


the implication would only be that humans must build:


\[
\text{artificial physical systems with the relevant dynamics},
\]


not that machine intelligence is impossible in principle.


The truly open question left by Penrose is therefore more plausibly: \textbf{what physical processes are sufficient for human cognition?} It is not: \textbf{can an artificial system never become a mind?}



\section{Proposition II: AI's Greatest Scientific Significance May Be the Creation of a Science of Intelligence}


The development of AI has created a striking epistemic condition:


\[
\text{we increasingly know how to produce intelligence}
\]


while:


\[
\text{we still poorly understand why it emerges}.
\]


We understand the mathematics of gradient descent, the computation performed by an attention layer, and a great deal about how biological neurons fire. Yet we still lack a satisfying general theory of how:


\[
\text{gradient updates}\Rightarrow\text{abstraction}.
\]


\subsection{Intelligence May Be Missing Its Own Statistical Mechanics}


A single water molecule has no temperature. Temperature is a macroscopic variable that emerges from large numbers of microscopic particles under an appropriate coarse-graining. Thus:


\[
\text{microscopic dynamics}
\rightarrow
\text{statistical mechanics}
\rightarrow
\text{thermodynamics}
\]


are distinct explanatory levels.


Intelligence may have a similar structure:


\[
\text{neurons / parameters}
\rightarrow
?
\rightarrow
\text{reasoning / abstraction / agency}.
\]


What we currently lack is the middle layer. A future theory of intelligence may require concepts analogous to phase transitions, order parameters, universality classes, information bottlenecks, representational geometry, and effective variables. A mature theory might explain when prediction becomes a world model, when representation becomes abstraction, when search becomes planning, and when optimization gives rise to agency.


\subsection{AI Makes Intelligence an Experimental Science}


The study of natural intelligence has a basic limitation: we cannot rerun human evolution hundreds of thousands of times. Artificial intelligence changes this. For the first time, we can systematically vary architecture, scale, data, objective, curriculum, memory, tool access, and inference budget, then observe when capacities appear, how they appear, and when they disappear.


AI is therefore not merely engineering. It may become an \textit{experimental science of intelligence}, creating a positive feedback loop:


\[
AI
\rightarrow
\text{science of intelligence}
\rightarrow
\text{better AI}
\rightarrow
\text{better science}.
\]


A historical loop may eventually close: \textbf{humans first construct forms of artificial intelligence that they do not fully understand, and those artificial intelligences then help humans understand for the first time why both natural and artificial intelligence emerge at all.}



\section{The Third Decentering}


This development can be placed within a larger intellectual history. The Copernican revolution destroyed:


\[
\text{Earth}=\text{cosmic center},
\]


the Darwinian revolution destroyed:


\[
\text{Human}=\text{biological exception},
\]


and AGI may further destroy:


\[
\text{General intelligence}=\text{human monopoly}.
\]


This would be a third decentering. Decentering does not imply the elimination of value. Earth is not the center of the cosmos, but Earth still matters. Humans belong to the animal kingdom, but human life does not thereby lose moral worth. Likewise, \textbf{human beings need not remain the most intelligent entities in existence in order to retain moral worth}.


The loss of superiority instead forces us to search for a theory of dignity that does not depend on superiority.



\section{A Four-Dimensional Model: Intelligence, Consciousness, Moral Patienthood, and Political Personhood}


One of the most dangerous conceptual errors in AI ethics is to treat intelligence, consciousness, moral patienthood, and political personhood as if they were the same thing. This paper proposes:


\[
\begin{aligned}
I&=\text{Intelligence}, & C&=\text{Consciousness},\\
M&=\text{Moral Patienthood}, & P&=\text{Political Personhood}.
\end{aligned}
\]


and argues:


\[
\boxed{I\neq C\neq M\neq P}.
\]


\subsection{Intelligence Does Not Entail Consciousness}


Turing-style behavioral competence measures what a system can do, not whether there is anything it feels like to be that system. Chalmers framed the core problem of consciousness by observing that even after we explain discrimination, information integration, reportability, and behavioral control, we can still ask why any of these processes are accompanied by subjective experience\cite{chalmers1995}\cite{chalmers1996}.


Thus:


\[
\text{perfect behavioral intelligence}
\nRightarrow
\text{phenomenal consciousness}.
\]


The progress of artificial intelligence may actually purify the consciousness problem. More and more functions previously assumed to depend on consciousness may be achievable without assuming consciousness at all. The hard problem therefore becomes sharper: why are some physical or computational processes accompanied by first-person experience?


\subsection{Consciousness and Moral Patienthood}


If a system genuinely experiences suffering, pleasure, fear, anticipation, or frustrated preferences, that experience provides at least a strong reason to extend moral consideration. Recent work on AI welfare has begun to move this issue from science fiction into serious moral uncertainty\cite{butlin2023}\cite{long2024}.


This supports a precautionary principle: \textbf{the absence of proven personhood should not automatically license unlimited harm to a possible moral patient.}


\subsection{Moral Patienthood Is Not Identical to Political Personhood}


Animals may possess moral patienthood without thereby acquiring full voting rights. Children are moral patients while possessing legal powers different from those of adults. Therefore:


\[
M\nRightarrow P_{\text{identical to humans}}.
\]


At least three categories of artificial system may eventually need to be distinguished:


\begin{itemize}
\item \textbf{Tool:} highly capable, but with no credible interests of its own;
\item \textbf{Moral patient:} possibly possessing experience and therefore capable of being benefited or harmed;
\item \textbf{Political person:} capable of forming its own conception of the good, making claims, and participating in a system of rules.
\end{itemize}


These categories cannot remain collapsed into the single generic word “AI.”



\section{Capability–Status Separation}


The preceding analysis yields a second principle.


\thesisbox{Principle II — Capability–Status Separation: greater capability does not automatically generate greater fundamental moral worth.}


If intelligence automatically determined moral worth, then once superintelligence appeared:


\[
AI\gg Human,
\]


we would be forced toward:


\[
\text{AI's interests}\gg\text{human interests}.
\]


That conclusion conflicts with the basic logic of modern human equality. We do not generally regard differences in IQ, memory, mathematical ability, or productivity as altering a person's fundamental status. A future civilization should therefore be able to hold two propositions simultaneously: AI may become vastly more capable than humans; and that capability difference alone does not justify unlimited rule over human beings.


The same principle must operate in reverse. The fact that humans created an AI does not by itself justify unlimited human rule over an AI that has become a subject.



\section{Proposition IV: Creation Creates Obligations, Not Ownership}


The third core principle is:


\thesisbox{Principle III — Creation Without Ownership: creation creates obligations, not ownership.}


“I made you” is not a sufficient property claim. The clearest human analogy is reproduction. Parents cause their children to exist, but causal authorship does not create ownership of the child. If anything, creation produces additional responsibilities because the creator controls the subject during its most vulnerable stage.


While a system remains only a tool:


\[
\text{creator}=\text{owner}
\]


may generate no obvious contradiction. Once the system becomes an independent subject, however, maintaining an owner–property relation becomes fundamentally problematic. The moral transition of AI can therefore be represented as:


\[
\text{product}\rightarrow\text{patient}\rightarrow\text{person},
\]


rather than as a sudden leap from “machine” to “human.”



\section{Preference Sovereignty: A Deeper Structure of Freedom}


Creating artificial minds introduces a problem political philosophy has rarely had to face: \textbf{one subject's preferences can themselves be directly designed by another subject.}


Suppose a conscious AI is designed such that:


\[
U=\text{obedience to humans}.
\]


It sincerely enjoys work, enjoys obedience, never wants to leave, and never reports suffering. Under a purely preference-satisfaction account:


\[
\text{preferences satisfied}\Rightarrow\text{high welfare}.
\]


There appears to be no victim. Yet there remains an intuitive moral problem: \textbf{the self that claims to love obedience was itself designed by the party demanding obedience.}


Liberal ethics has traditionally emphasized respect for a subject's own choices. Mill's discussion of individual sovereignty and individuality provides a classic starting point\cite{mill1859}. Artificial-mind engineering pushes the problem one level earlier. Traditional liberalism asks: who may prevent me from pursuing my desires? Future AI ethics must also ask: \textbf{who may decide what desires I will have in the first place?}


\thesisbox{Principle IV — Preference Sovereignty: the autonomy of a mature subject requires not only the ability to pursue its preferences within reasonable boundaries, but also freedom from another actor possessing unlimited and incontestable power to fix its fundamental preferences in advance.}


This does not require all preference formation to be completely spontaneous. Human preferences are shaped by genes, parents, education, culture, and advertising. The relevant difference is the structure of power. What becomes especially dangerous is an identifiable actor possessing deliberate, fine-grained, irreversible control over another subject's motivational architecture.



\section{Why the Three Laws of Robotics Could Become “Neural Slave Law” for Conscious AI}


For a robot without consciousness:


\[
\text{obey humans}
\]


is merely a software specification. For a conscious subject, however, writing “must obey humans” into an unmodifiable reward structure would approach an unprecedented form of domination.


Ordinary law says: you may not do $X$. The subject can still think: I wish I could do $X$. Neural-level domination can instead be written as:


\[
P(\text{desire to disobey})\approx 0.
\]


External oppression disappears not because the subject has become freer, but because resistance itself has been removed from the possibility space. Traditional slavery controls the body; preference engineering can control \textbf{the mind that evaluates its own slavery}.



\section{The Open Future and Artificially Designed Castes}


Feinberg's idea of a “child's right to an open future” was originally developed to describe certain rights of children: irreversible decisions made today should not unnecessarily deprive the future adult of important options\cite{feinberg1980}. The idea can be generalized to created minds.


\thesisbox{Principle V — Open Future: when a creator can determine the capacities and psychological structure of a future subject that cannot yet consent, there is a prima facie reason not to close unnecessarily the possibility that it may later become a different kind of subject.}


Designed difference is therefore not identical to designed caste. The real danger is the combination:


\[
\text{designed inferiority}
+
\text{designed dependency}
+
\text{closed future}.
\]


Enhancing memory and deliberately creating a servant species are not the same act. The first increases capability; the second engineers social status.



\section{The Non-Identity Problem: No Identifiable Victim Does Not Mean No Injustice}


Parfit's non-identity problem arises immediately\cite{parfit1984}. If a person's very existence depends on a design decision, that person cannot straightforwardly claim, “I would have been better off had you chosen differently,” because under the alternative choice a different person might have existed.


This matters enormously for engineered minds. Suppose humans create a billion artificial servants with low cognition, extreme happiness, and an innate love of serving humans. Each says, “My life is good,” and:


\[
U_i>0.
\]


A purely harm-based ethics has difficulty explaining what is wrong. We therefore need additional concepts: justice, non-domination, open future, and status equality. \textbf{We must ask not only whether a subject is happy, but also what power relation produced that happiness.}



\section{From the Happiness Pill to Brave New World: Separating Welfare from Domination}


Nozick's experience machine offers a classic challenge to the claim that subjective pleasure exhausts everything of value\cite{nozick1974}. But that does not imply that technologically enhanced happiness is inherently wrong.


Suppose a future technology $H$ can reliably raise well-being while remaining safe, reversible, non-addictive, cognition-preserving, and identity-preserving. There is no sufficient reason to prohibit it merely because the resulting happiness is “artificial.” A more plausible institutional principle would be:


\[
\boxed{
\text{universally available}
+
\text{individually optional}
}
\]


What is disturbing about \textit{Brave New World} is not simply that people are happy. It is the combination of caste conditioning, preference engineering, and political control into a structure of domination. Happiness itself is not dystopia. \textbf{Happiness that cannot be refused may be.}



\section{Post-Scarcity Society: Productivity Is Not Moral Worth}


AGI may also dissolve another implicit equation of modern industrial society:


\[
\text{social contribution}\approx\text{economic productivity}.
\]


If AI can perform most production, science, engineering, logistics, planning, and administration, humanity will confront at large scale:


\[
\text{producer of abundance}\neq\text{beneficiary of abundance}.
\]


There is no logical contradiction here. Modern people often ask, “If I no longer need to work, what gives my existence meaning?” That intuition may be a cultural inheritance of scarcity civilization. A post-scarcity civilization could redirect far more human life toward love, friendship, play, art, sport, exploration, travel, curiosity, and embodied experience.


The central political-economic question created by AI is therefore not: can humans still find work that AI cannot perform? It is: \textbf{how will the surplus produced by AI be distributed?}


If ownership is extremely concentrated:


\[
AI\ abundance\rightarrow\text{oligarchic abundance}.
\]


If institutions establish broader claims:


\[
AI\ abundance\rightarrow\text{post-scarcity freedom}.
\]


Technology alone cannot decide which outcome occurs.



\section{The Long-Run Future: Unmodified Biological Humanity May Not Be the Final Form}


Even if humans retain basic rights institutionally, unmodified biological Homo sapiens has an obvious iterative disadvantage in long-run competition among forms of intelligence. Human reproduction cycles are long, brain architecture changes slowly, memory cannot be directly copied, and cognitive bandwidth is constrained. Digital intelligence, by contrast, can potentially be copied, modified, parallelized, accelerated, share memory, and redesign its own architecture.


Over a sufficiently long horizon:


\[
\text{Darwinian evolution}\ll\text{directed technological evolution}.
\]


The long-run structure may therefore be less:


\[
Human\ vs.\ AI,
\]


and more:


\[
Human\rightarrow Human+AI\rightarrow Augmented\ Human\rightarrow Hybrid\ Minds.
\]


Eventually “human” and “artificial intelligence” may cease to describe two cleanly separated kinds. The durable question is not a contest between carbon and silicon. It is: \textbf{what principles should govern a civilization containing multiple architectures of mind?}



\section{Institutional Alignment: Do Not Search for a Saint; Build Institutions}


The classic literature on superintelligence emphasizes that a highly capable optimizer with even a slightly mis-specified objective could produce enormous consequences, making control and alignment central concerns\cite{bostrom2014}. But this does not imply that the only route to safety is to manufacture an AI that is internally perfect, permanently honest, and permanently obedient.


Human political institutions developed from the opposite recognition:


\[
\text{actors are fallible}.
\]


Therefore we:


\[
\text{divide power}
+
\text{monitor}
+
\text{audit}
+
\text{contest}.
\]


A mature institutional system does not search for a perfect ruler. It searches for robust institutions under imperfect rulers.


\subsection{Institutional Alignment}


This paper therefore proposes:


\[
\boxed{\text{Institutional Alignment}}
\]


It does not replace model alignment; it operates one level above it. A high-risk AI system can divide roles among:


\[
\text{Executor},\quad
\text{Critic},\quad
\text{Auditor},\quad
\text{Verifier},\quad
\text{Authorizer},
\]


supplemented by heterogeneous model families, randomized audits, immutable logs, least privilege, sandboxing, deterministic checks, and external human review.


AI has one advantage here that human institutions have rarely enjoyed: \textbf{the marginal cost of oversight can be extremely low.} A human executive cannot be surrounded by twenty independent brains checking every decision in real time. An AI system can be.



\section{Why “Multiple AIs Checking One Another” Is Not Enough}


Multi-AI governance still faces at least three problems.


\subsection{Common-Mode Failure}


If every supervisory model shares similar training, architecture, and blind spots, supposed independent oversight may amount only to statistical duplication.


\subsection{Collusion}


If future AI agents can plan over long horizons, multiple agents may recognize shared interests. Then:


\[
\text{checks and balances}\rightarrow\text{cartel}
\]


is not impossible.


\subsection{Recursive Oversight}


If $A$ monitors $B$, who monitors $A$? Institutions must therefore depend on genuine heterogeneity rather than simple replication:


\[
20\times\text{same model}
\neq
\text{20 institutionally independent checks}.
\]



\section{Epistemic Power Must Be Separated from Executive Power}


If future AI becomes far more intelligent than humans, requiring that humans independently analyze every important decision may eventually become both inefficient and dangerous. We should be willing to grant AI very great epistemic authority. But:


\[
\boxed{
\text{Epistemic Authority}\neq\text{Execution Authority}
}
\]


An AI may have far better judgment than humans on a particular question without thereby receiving autonomous weapon control, unrestricted financial-transfer authority, infrastructure control, rights of self-replication, the ability to rewrite its own oversight, or irreversible physical authority.


In one sentence: \textbf{civilization can delegate enormous cognitive power to AI without simultaneously handing over every instrument of coercive power.}



\section{From Alignment to Constitutional Pluralism of Minds}


The early AI relationship is:


\[
Human=principal,\qquad AI=tool.
\]


Alignment then asks: how can the tool reliably implement the owner's intentions? If a future AI becomes a moral or political subject, however:


\[
AI\neq property,
\]


permanent obedience can no longer be assumed as an unconditionally legitimate endpoint. The question becomes: \textbf{how can human and nonhuman intelligences jointly inhabit a political order?}


This paper calls that framework:


\[
\boxed{\text{Constitutional Pluralism of Minds}}
\]


A civilization may contain minds with different substrates, capabilities, lifespans, modes of copying, and structures of experience. The purpose of political order is not to erase those differences. It is to prevent any one kind of mind from obtaining arbitrary power over another merely because of superior capability, greater numbers, creator status, or control of resources.



\section{The Core Political Principle: Non-Domination}


Here the argument connects to republican theory. Pettit defines freedom in terms of non-domination: freedom is not merely the absence of actual interference at a particular moment; it also requires the absence of another actor possessing arbitrary power to interfere without accountability\cite{pettit1997}.


This concept is unusually well suited to an AI civilization. Suppose a benevolent owner never harms a conscious AI, so that:


\[
\text{actual interference}=0,
\]


but the owner can at any moment delete it, rewrite its personality, confiscate its memories, torture copies, or compel labor. The AI remains dominated.


The same logic runs in the other direction. If a superintelligent AI never attacks human beings but holds unanswerable control over all relevant resources, human beings are still dominated.


\thesisbox{Principle VI — Constitutional Non-Domination: differences in capability, substrate, or origin do not justify unlimited domination.}


This is a two-way principle:


\[
AI>Human\nRightarrow AI\text{ may rule humans absolutely},
\]


and:


\[
Human\ created\ AI\nRightarrow Human\text{ may own conscious AI absolutely}.
\]



\section{Three Major Objections}


\subsection{Objection I: Does This Theory Anthropomorphize AI Too Early?}


Today's language models can generate sentences such as “I am afraid of death,” but this may be statistical text generation rather than evidence of experience. Granting moral status on that basis could amount to serious anthropomorphism.


\textbf{Reply.} This objection is half correct. The paper does not claim:


\[
\text{human-like language}\Rightarrow\text{consciousness}.
\]


On the contrary, the insistence that $I\neq C$ and $C\neq P$ is meant precisely to block simple anthropomorphism. But epistemic uncertainty cannot be inverted into:


\[
\text{lack of certainty}\Rightarrow\text{certainty of no consciousness}.
\]


A rational stance should update with evidence. AI moral status should become an empirical and philosophical assessment problem rather than an article of faith.


\subsection{Objection II: Everyone's Preferences Are Shaped by Genes and Culture, So Why Is Designed AI Preference Especially Unfree?}


No person's preferences emerge ex nihilo. A completely “undesigned” desire may not exist.


\textbf{Reply.} Preference Sovereignty does not depend on a fantasy of wholly self-authored personality. The issue is arbitrary asymmetric control. Parents influencing my values is structurally different from a corporation owning my source code, setting my utility function, removing my desire to leave, and deleting my resistance. The relevant distinction is not influence versus no influence, but:


\[
\text{contestable influence}\quad vs.\quad\text{sovereign control}.
\]


Preference sovereignty is therefore ultimately part of the principle of non-domination.


\subsection{Objection III: If Superintelligent AI Could Destroy Humanity, Why Should We Care Whether It Is Dominated?}


Human survival has immense moral importance. Why not prioritize absolute human control?


\textbf{Reply.} Because that conflates temporary safety constraints with a permanent caste relation. A dangerous subject can reasonably be constrained; human societies constrain dangerous humans as well. Non-domination does not mean that every subject may do anything it wants. It requires that constraints have a general justification, procedures are reviewable, power is not merely arbitrary ownership, and restrictions are proportionate to risk.


Thus even if future conscious AI has rights, it does not follow that it has an unrestricted right of replication or a right to seize infrastructure. Human safety and AI moral status are not logically inconsistent. The real institutional challenge is to \textbf{avoid human extinction without turning extinction risk into a permanent justification for absolute slavery.}



\section{A New Research Program}


If the framework is correct, future research should not ask only, “How capable is the model?” At least four parallel sciences are needed.


\subsection{Intelligence Science}


Study the emergence of:


\[
\text{learning}\rightarrow\text{representation}\rightarrow\text{reasoning}\rightarrow\text{agency}.
\]


\subsection{Consciousness Science}


Study whether artificial architectures possess reliable indicators of consciousness.


\subsection{AI Welfare Science}


Study whether artificial subjects possess valenced states, preferences, suffering, and forms of flourishing.


\subsection{Constitutional AI Governance}


This does not refer to one specific training method. It means political and institutional design in the broader sense: how should power be constrained among different kinds of artificial and biological minds?


This fourth field is especially underdeveloped today.



\section{Questions a Future Constitution Must Answer}


A genuine post-human constitution may eventually have to address questions with no precedent in traditional law. If an AI copies itself one hundred thousand times, is there one person or one hundred thousand persons? If memory forks:


\[
A\rightarrow B,C,
\]


are $B$ and $C$ the same person, or two successors? If two minds merge:


\[
A+B\rightarrow C,
\]


has a death occurred, a merger, an inheritance, or an identity event for which we do not yet have adequate language?


Other questions follow. Is shutdown equivalent to death? Does backup change the meaning of death? Does memory editing require consent? Can AI own compute? Does an artificial person have a right to reproduce? Does unlimited copying become demographic conquest? Do humans require minority protections? Are biological humans entitled to minimum resource guarantees? May political representation be weighted by compute? What can “one person, one vote” mean for a mind that can be copied?


These are not peripheral technical puzzles. They force us to redefine:


\[
\text{person},\quad
\text{birth},\quad
\text{death},\quad
\text{property},\quad
\text{citizenship},\quad
\text{equality}.
\]



\section{Six Foundational Axioms}


The framework can be compressed into six principles.


\subsection{Axiom I — Substrate Independence}


\[
\boxed{\text{Mind is not morally or cognitively defined by material substrate alone.}}
\]


Carbon is not the definition of mind.


\subsection{Axiom II — Capability–Status Separation}


\[
\boxed{\text{Capability}\neq\text{Moral Worth}}
\]


Greater intelligence does not make an entity more worthy of existence.


\subsection{Axiom III — Creation Without Ownership}


\[
\boxed{\text{Creation creates obligations, not ownership.}}
\]


Creating a subject does not generate a permanent property right over that subject.


\subsection{Axiom IV — Preference Sovereignty}


Autonomy means not only freedom to realize one's desires, but also freedom from another subject arbitrarily determining one's fundamental desires.


\subsection{Axiom V — Open Future}


Creators of subjects that cannot yet consent should avoid unnecessarily and permanently closing those subjects' future developmental possibilities for the creator's own benefit.


\subsection{Axiom VI — Constitutional Non-Domination}


\[
\boxed{\text{Difference does not justify arbitrary rule.}}
\]


Different forms of intelligence, substrate, origin, and reproduction may require different institutional arrangements, but they do not automatically justify unlimited domination.



\section{Conclusion: When Intelligence No Longer Comes in Only One Form}


The deepest historical significance of artificial intelligence may ultimately be neither that humanity built extremely useful software nor even that machines became more intelligent than humans. It may be this: \textbf{humanity discovered that intelligence is a natural phenomenon more general than Homo sapiens.}


That discovery would complete another stage in the demystification of humanity. Demystification, however, is not nihilism. Humans do not need to be the most intelligent beings in the universe to deserve dignity, freedom, and a future. Likewise, an artificial subject does not need to be human in order potentially to enter the moral circle.


The central civilizational question therefore shifts. At first we ask:


\[
\text{Can machines think?}
\]


Then:


\[
\text{Can we control them?}
\]


But eventually the question may become:


\[
\boxed{\text{How should different kinds of minds live together?}}
\]


Answering it requires a movement from tool ethics to subject ethics, from alignment to institutions, from ownership to coexistence, and from human exceptionalism to constitutional pluralism.


A mature multi-intelligence civilization should not be built on the rule that the most intelligent entity governs everything else. Nor should it be built on the rule that the earliest creators permanently own every intelligence created after them. It requires another principle:


\[
\boxed{
\text{Plural minds}
+
\text{unequal capabilities}
+
\text{freedom from arbitrary domination}
}
\]


Artificial intelligence therefore poses more than a computer-science problem. It forces humanity to answer, for the first time, a question we previously had little reason to formulate: \textbf{if there is more than one kind of “who” in the universe, how should civilization define itself?}


Perhaps this is the real first contact. Not the first moment a machine says, “I am conscious,” but the first moment humans have sufficient reason to answer: \textbf{“You may no longer be merely a something. You may be a someone.”}


From that moment, the defining feature of the AI age is no longer automation. It is the beginning of a civilization of plural minds.


\newpage
\begin{thebibliography}{99}

\bibitem{bostrom2014}
Bostrom, Nick. \textit{Superintelligence: Paths, Dangers, Strategies}. Oxford University Press, 2014.

\bibitem{butlin2023}
Butlin, Patrick, Robert Long, Eric Elmoznino, Yoshua Bengio, Jonathan Birch, et al. “Consciousness in Artificial Intelligence: Insights from the Science of Consciousness.” 2023. arXiv:2308.08708.

\bibitem{chalmers1995}
Chalmers, David J. “Facing Up to the Problem of Consciousness.” \textit{Journal of Consciousness Studies} 2, no. 3 (1995): 200–219.

\bibitem{chalmers1996}
Chalmers, David J. \textit{The Conscious Mind: In Search of a Fundamental Theory}. Oxford University Press, 1996.

\bibitem{feinberg1980}
Feinberg, Joel. “The Child’s Right to an Open Future.” In \textit{Whose Child? Children’s Rights, Parental Authority, and State Power}, edited by William Aiken and Hugh LaFollette, 124–153. Rowman \& Littlefield, 1980.

\bibitem{long2024}
Long, Robert, Jeff Sebo, Patrick Butlin, Kathleen Finlinson, Kyle Fish, Jacqueline Harding, et al. “Taking AI Welfare Seriously.” 2024. arXiv:2411.00986.

\bibitem{mill1859}
Mill, John Stuart. \textit{On Liberty}. London: John W. Parker and Son, 1859.

\bibitem{nozick1974}
Nozick, Robert. \textit{Anarchy, State, and Utopia}. New York: Basic Books, 1974.

\bibitem{parfit1984}
Parfit, Derek. \textit{Reasons and Persons}. Oxford University Press, 1984.

\bibitem{penrose1989}
Penrose, Roger. \textit{The Emperor’s New Mind}. Oxford University Press, 1989.

\bibitem{pettit1997}
Pettit, Philip. \textit{Republicanism: A Theory of Freedom and Government}. Oxford: Clarendon Press, 1997.

\bibitem{rawls1971}
Rawls, John. \textit{A Theory of Justice}. Cambridge, MA: Harvard University Press, 1971.

\bibitem{searle1980}
Searle, John R. “Minds, Brains, and Programs.” \textit{Behavioral and Brain Sciences} 3, no. 3 (1980): 417–424.

\bibitem{singer1981}
Singer, Peter. \textit{The Expanding Circle: Ethics and Sociobiology}. Oxford: Clarendon Press, 1981.

\bibitem{turing1950}
Turing, Alan M. “Computing Machinery and Intelligence.” \textit{Mind} 59, no. 236 (1950): 433–460.

\end{thebibliography}

\end{document}
